\chapter{Atmosphere}
\label{cap:atmosfera}

\section{Pressure-level variables}
\label{sec:livelli-pressione}

\Cref{tab:var-pressione} lists the 13 pressure-level variables, 169 messages in all. Besides
temperature, wind, humidity and geopotential, they include the two forms of vertical velocity, in
pressure (\campo{w}) and geometric (\campo{wz}), the cloud fraction and four cloud and
precipitation water contents.

\begin{table}[htbp]
  \centering
  \small
  \caption{Variables on the 13 pressure levels.}
  \label{tab:var-pressione}
  \begin{tabular}{@{}lrll@{}}
    \toprule
    Field & paramId & Description & Units \\
    \midrule
    \campo{cc}   & 248    & cloud fraction & 0--1 \\
    \campo{ciwc} & 247    & specific cloud ice water content & \si{kg/kg} \\
    \campo{clwc} & 246    & specific cloud liquid water content & \si{kg/kg} \\
    \campo{crwc} & 75     & specific rain water content & \si{kg/kg} \\
    \campo{cswc} & 76     & specific snow water content & \si{kg/kg} \\
    \campo{q}    & 133    & specific humidity & \si{kg/kg} \\
    \campo{r}    & 157    & relative humidity & \si{\percent} \\
    \campo{t}    & 130    & temperature & \si{K} \\
    \campo{u}    & 131    & zonal wind component & \si{m/s} \\
    \campo{v}    & 132    & meridional wind component & \si{m/s} \\
    \campo{w}    & 135    & vertical velocity in pressure & \si{Pa/s} \\
    \campo{wz}   & 260238 & geometric vertical velocity & \si{m/s} \\
    \campo{z}    & 129    & geopotential & \si{m^2/s^2} \\
    \bottomrule
  \end{tabular}
\end{table}

\paragraph{Specific contents.} WRF carries vapor and hydrometeors as \emph{mixing ratios}, per
kilogram of dry air. ECMWF defines \campo{q} and the water contents as \emph{specific contents},
per kilogram of moist air, and the mass of moist air comprises dry air, vapor, cloud water, ice,
rain and snow \parencite{paramdb}. The archive therefore divides \campo{q}, \campo{clwc},
\campo{ciwc}, \campo{crwc}, \campo{cswc} and the column integrals by the same denominator,
$1 + r_t$, where $r_t$ is the total water mixing ratio. On the three-dimensional field the
correction is \SI{0.011}{\percent} in the median, \SI{1.44}{\percent} at the 99th percentile and
\SI{2.26}{\percent} at most.

\paragraph{Cloud fraction.} With the Xu--Randall scheme the cloud fraction \campo{cc} is
continuous. On a subdomain of $500 \times 600$ cells only between \SI{18}{\percent} and
\SI{22}{\percent} of the cloudy points are exactly 1, and the others span the whole interval from
0 to 1.

\paragraph{Relative humidity.} The field \campo{r} is ECMWF's relative humidity, computed over the
mixed phase with equations 7.5, 7.6, 7.89 and 7.90 of the IFS Cy41r2 documentation, the cycle
with which ERA5 was produced \parencite{era5doc}, and with the constants of that cycle
\parencite{ifs41r2p4}. The saturation vapor pressure is
\begin{equation}
  e_\text{sat}(T) = \alpha(T)\, e_\text{sat,w}(T) + \bigl(1 - \alpha(T)\bigr)\, e_\text{sat,i}(T),
  \qquad
  \alpha(T) = \left(\frac{T - T_\text{ice}}{T_0 - T_\text{ice}}\right)^{2},
  \label{eq:fase-mista}
\end{equation}
where $T$ is the temperature, $e_\text{sat,w}$ and $e_\text{sat,i}$ are the saturation vapor
pressures over water and over ice (Tetens formula), $T_\text{ice} = \SI{250.16}{K}$,
$T_0 = \SI{273.16}{K}$, and $\alpha$ is 0 below $T_\text{ice}$ and 1 above $T_0$. Below
$T_\text{ice}$ this relative humidity differs by up to a factor of \num{1.8} from the one computed
over liquid water. The field is not truncated at \SI{100}{\percent}, because supersaturation is a
state of the model and ERA5 publishes it too. In the hourly file of 15 July 2024 at 12 UTC,
\campo{r} reaches \SI{139.7}{\percent} at \SI{300}{\hPa} in the ice regime, and \num{34785} points
exceed \SI{100}{\percent} at \SI{400}{\hPa}. An ingestion pipeline that discards values above
\SI{100}{\percent} discards valid data.

\paragraph{Geopotential.} The published geopotential is $z = g_0\, z_\text{WRF}$, with
$g_0 = \SI{9.80665}{m/s^2}$ the standard gravity and $z_\text{WRF}$ the geopotential height of the
model. Dividing \campo{z} by $g_0$ recovers the model height exactly, at the surface and on every
pressure level. The surface geopotential and the pressure-level geopotential follow the same
convention, verified point by point on the \num{1349910} land cells with elevation above \SI{1}{m} on
15 July 2024 at 12 UTC. WRF, however, integrates the hydrostatic equation with
$g = \SI{9.81}{m/s^2}$, so the archive departs by \SI{0.034}{\percent} from hydrostatic balance
with its own \campo{t} and pressure. This is one part in \num{2900}, about \SI{2}{m} at
\SI{500}{\hPa} and \SI{7}{m} at \SI{50}{\hPa} \parencite{issue39}. The two properties, \campo{z}
divided by $g_0$ equal to the model height and hydrostatic balance with the model's \campo{t} and
pressure, cannot hold together.

\section{Below ground}
\label{sec:sotto-terreno}

Where a pressure level lies below ground, the archive publishes the value of the lowest model
level as it is, without extrapolating or filling (\emph{clamping}). Below ground \campo{t}
coincides exactly with the lowest model level, \campo{q} within \SI{1e-9}{kg/kg} and \campo{z}
within \SI{2.5e-4}{m} \parencite{issue44}. The value is a real state of the model, taken between
\SI{24}{m} and \SI{27}{m} above ground, and it stays consistent with the near-surface fields. In
July, $\campo{2t} - \campo{t}(\SI{1000}{\hPa})$ below ground is \SI{-0.9}{K} at midnight and
\SI{+1.4}{K} at noon, that is, a nocturnal inversion and a daytime superadiabatic layer in the
roughly \SI{22}{m} between the \SI{2}{m} level and the lowest model level. Below ground \campo{r}
reaches about \SI{163}{\percent}, because it pairs the humidity of the lowest model level with the
nominal pressure of \SI{1000}{\hPa}.

The cost is that \campo{z} contradicts \campo{sp}. Where the surface pressure places
\SI{1000}{\hPa} below ground, the geopotential published at \SI{1000}{\hPa} places it, in the
median, between \SI{24}{m} and \SI{27}{m} above. This occurs on \SI{47.6}{\percent} of the domain
at \SI{1000}{\hPa} and on \SI{9}{\percent} at \SI{925}{\hPa}. ERA5 instead extrapolates below
ground with a temperature lapse rate, and the two archives differ by convention precisely where
\SI{1000}{\hPa} is below ground. The whole difference at \SI{1000}{\hPa} comes from this: on the
verification sample the bias of \campo{t} against ERA5 is \SI{-1.71}{K} at \SI{1000}{\hPa} against
\SI{-0.05}{K} at \SI{850}{\hPa}, and the correlation of \campo{z} with ERA5 at \SI{1000}{\hPa} is
\num{0.108}. For a comparison with ERA5 or for computing thicknesses, the below-ground levels
must be masked with \campo{sp}, which is in the same hourly file.

\section{The \SI{50}{\hPa} level}
\label{sec:livello-50}

The model top is at \SI{50}{\hPa} (\cref{sec:dominio}), and the highest model level at which WRF
defines temperature always lies below it, for example at \SI{51.9}{\hPa} on average over the domain
on 19 January 2024 at 00 UTC. The \SI{50}{\hPa} pressure level therefore lies above the highest
model level, and the vertical interpolation repeats the
values of that level there, as the clamping does below ground (\cref{sec:sotto-terreno}). At that
validity time \campo{t} at \SI{50}{\hPa} equals the temperature of the highest model level to
within the encoding precision, \SI{0.009}{K}, and $\campo{z}/g_0$ equals its geopotential height to
within \SI{1}{m}; the same holds for the other fields at \SI{50}{\hPa}. At the same validity time the
layer between the highest model level and \SI{50}{\hPa} has a hydrostatic thickness of \SI{227}{m}
on average, between \SI{214}{m} and \SI{239}{m} over the domain, and $\campo{z}/g_0$ at
\SI{50}{\hPa} lies that much below the height of the \SI{50}{\hPa} pressure level. The thickness
between \SI{100}{\hPa} and \SI{50}{\hPa} computed from the archive, \SI{4009}{m} on average, is
therefore about \SI{5.4}{\percent} short of the hydrostatic thickness, about \SI{4236}{m}. A
geostrophic wind computed from \campo{z} at \SI{50}{\hPa} is affected by the horizontal variation of
the missing thickness, \SI{25}{m} between its minimum and maximum over the domain.

\section{Near-surface fields}
\label{sec:vicino-superficie}

\Cref{tab:var-quota} lists the fields at fixed height above ground and the mean sea level
pressure.

\begin{table}[htbp]
  \centering
  \small
  \caption{Fields at fixed height and mean sea level pressure.}
  \label{tab:var-quota}
  \begin{tabular}{@{}lrlll@{}}
    \toprule
    Field & paramId & Description & Units & Height \\
    \midrule
    \campo{10u}, \campo{10v}   & 165, 166       & wind components & \si{m/s} & \SI{10}{m} \\
    \campo{100u}, \campo{100v} & 228246, 228247 & wind components & \si{m/s} & \SI{100}{m} \\
    \campo{200u}, \campo{200v} & 228239, 228240 & wind components & \si{m/s} & \SI{200}{m} \\
    \campo{10fg} & 49     & hourly maximum of the resolved wind & \si{m/s} & \SI{10}{m} \\
    \campo{2t}   & 167    & temperature & \si{K} & \SI{2}{m} \\
    \campo{2d}   & 168    & dewpoint temperature & \si{K} & \SI{2}{m} \\
    \campo{2r}   & 260242 & relative humidity & \si{\percent} & \SI{2}{m} \\
    \campo{vwsh} & 260068 & vertical wind shear between \SI{10}{m} and \SI{100}{m} & \si{s^{-1}} & \SI{100}{m} \\
    \campo{msl}  & 151    & mean sea level pressure & \si{Pa} & sea level \\
    \bottomrule
  \end{tabular}
\end{table}

\paragraph{Relative humidity at \SI{2}{m}.} \campo{2r} is the relative humidity over liquid water
at every temperature, the World Meteorological Organization convention for screen-level humidity
measurements \parencite{wmo8}, that is, for the observations against which a \SI{2}{m} humidity is
compared. ERA5 does not publish a relative humidity at \SI{2}{m}. \campo{r} and \campo{2r}
therefore follow two different definitions and differ where the temperature is below $T_0$. A
profile that places \campo{2r} at the base of \campo{r}, or a vertical humidity gradient computed
between the two, mixes two quantities. Recomputing the relative humidity from
\campo{q}, \campo{t} and \campo{sp}, all of which are published, puts both in the same convention.

\paragraph{Dewpoint temperature.} The inversion that derives the dewpoint temperature from the
vapor gives, at saturated points, a \campo{2d} that exceeds \campo{2t} by \SI{0.0133}{K}, the same
at every validity time and in every season. At four validity times in January 2024, July 2024
and February 2025 these points number between \num{6161} and \num{64005}, all with a relative
humidity of \SI{100}{\percent} \parencite{issue42}. The value is left as it is. Truncating
\campo{2d} at \campo{2t} would remove the visible crossings but leave the same offset at every
unsaturated point, introduce a kink in the field at \SI{100}{\percent} and break the consistency
between \campo{2d} and the model's \SI{2}{m} vapor (\campo{Q2}), which holds within a relative
error between \num{3.6e-4} and \num{1.1e-2} on the vapor pressure. A relative humidity
reconstructed from \campo{2t} and \campo{2d} therefore reaches \SI{100.09}{\percent}.

\section{Columns and clouds}
\label{sec:colonne}

\Cref{tab:var-colonna} lists the column integrals and the cloud covers.

\begin{table}[htbp]
  \centering
  \small
  \caption{Column integrals and cloud covers.}
  \label{tab:var-colonna}
  \begin{tabular}{@{}lrll@{}}
    \toprule
    Field & paramId & Description & Units \\
    \midrule
    \campo{tcwv} & 137    & total column water vapor & \si{kg/m^2} \\
    \campo{tclw} & 78     & total column cloud liquid water & \si{kg/m^2} \\
    \campo{tciw} & 79     & total column cloud ice water & \si{kg/m^2} \\
    \campo{tcrw} & 228089 & total column rain water & \si{kg/m^2} \\
    \campo{tcsw} & 228090 & total column snow water & \si{kg/m^2} \\
    \campo{tcw}  & 136    & total column water & \si{kg/m^2} \\
    \campo{tcc}  & 164    & total cloud cover & 0--1 \\
    \campo{hcc}  & 188    & high cloud cover & 0--1 \\
    \campo{mcc}  & 187    & medium cloud cover & 0--1 \\
    \campo{lcc}  & 186    & low cloud cover & 0--1 \\
    \bottomrule
  \end{tabular}
\end{table}

\paragraph{Total water.} ECMWF defines \campo{tcw} as the sum of five species, vapor, liquid
water, ice, rain and snow \parencite{paramdb}. The archive sums six species, because it includes
graupel \parencite{issue47}. ECMWF defines a graupel content (paramId 260028, \campo{grle}) and
its column (260001, \campo{tcolg}) \parencite{paramdb}, but ERA5 publishes neither, and graupel
appears only inside \campo{tcw}. Reconstructing \campo{tcw} as $\campo{tcwv} + \campo{tclw} +
\campo{tciw} + \campo{tcrw} + \campo{tcsw}$ does not reproduce it, and the residual is the graupel
column. At the four validity times examined, 15 January, April, July and October 2024 at 12 UTC, the
residual is, as a domain mean, \SI{0.4}{\percent} of the condensate. Its maximum is
\SI{5.69}{kg/m^2} in January, \SI{13.89}{kg/m^2} in April, \SI{16.73}{kg/m^2} in October and
\SI{26.46}{kg/m^2} in July, in a convective cell where it exceeds all the rest of the column's
condensate.

\paragraph{Cloud cover.} \campo{tcc}, \campo{hcc}, \campo{mcc} and \campo{lcc} are maximum-random
overlaps of WRF's cloud fraction \campo{CLDFRA}. On the verification sample, with the CHAPTER
fields aggregated onto the ERA5 grid, the ratio of means is \num{0.783} for total cloud cover and
\num{0.610} for high cloud cover (\cref{tab:nubi}), and CHAPTER total cloud cover is lower than ERA5's at
every one of the 34 validity times \parencite{issue48}.

\begin{table}[htbp]
  \centering
  \small
  \caption{Cloud covers and total column cloud ice water (\si{kg/m^2}): means over the
  verification sample, on the ERA5 grid.}
  \label{tab:nubi}
  \begin{tabular}{@{}lSSS@{}}
    \toprule
    Field & {CHAPTER} & {ERA5} & {Ratio of means} \\
    \midrule
    \campo{tcc}  & 0.381  & 0.487  & 0.783 \\
    \campo{hcc}  & 0.169  & 0.276  & 0.610 \\
    \campo{mcc}  & 0.161  & 0.205  & 0.786 \\
    \campo{lcc}  & 0.231  & 0.268  & 0.862 \\
    \campo{tciw} & 0.0185 & 0.0166 & 1.109 \\
    \bottomrule
  \end{tabular}
\end{table}

The gap does not come from the overlap assumption. Computed point by point on the archive's
\campo{CLDFRA}, the whole range of possible assumptions, from maximum to random overlap, is
\SIrange{2}{4}{\percent} wide, while ERA5 total cloud cover exceeds CHAPTER's by
\SI{28}{\percent}, so the assumption accounts for at most one seventh of the gap. The gap lies in
\campo{CLDFRA} itself, the Xu--Randall diagnostic. Total column cloud ice is \num{1.109} times
ERA5's while \campo{hcc} is \num{0.610} times, so the scheme assigns to the condensate present a
smaller cloud fraction than ERA5 does. Part of the gap comes from resolution. CHAPTER cloud cover is
overlapped at \SI{3}{km} and then aggregated, whereas ERA5 cloud cover is overlapped on a profile
already smooth at \SI{31}{km}, which for equal layer fractions gives a larger cloud cover. This effect
too is bounded by the same range, between \SI{2}{\percent} and \SI{4}{\percent}.

\section{Instability}
\label{sec:instabilita}

\campo{mucape} (paramId 228235) and \campo{mucin} (228236), in \si{J/kg}, are the convective
available potential energy (CAPE) and the convective inhibition of the most unstable parcel, that
is, the parcel mixed over \SI{500}{m} with the highest equivalent potential temperature in the
lowest \SI{3}{km}. They are not computed from the surface.

The routine that computes the two quantities returns an inhibition only in columns with CAPE of
at least \SI{100}{J/kg}, and elsewhere \campo{mucin} is a missing point (bitmap). A zero in those
columns would state that nothing suppresses convection precisely where it is suppressed most,
because a column without CAPE is normally a column with strong inhibition. Where \campo{mucin} is
missing, ERA5 carries a mean inhibition between \SI{25}{J/kg} and \SI{265}{J/kg}. Filling the
missing points with zeros would teach a statistical model that situations with little CAPE have
no inhibition, the opposite of the physics.

The fraction of the domain in which \campo{mucin} is present ranges from \SI{3.8}{\percent} in
January to \SI{28.8}{\percent} in July (\cref{tab:mucin}). Where it is present, \campo{mucin} is
between \num{0.03} and \num{0.24} times ERA5's, because the columns in which it exists are the
least inhibited.

\begin{table}[htbp]
  \centering
  \small
  \caption{Fraction of the domain in which \campo{mucin} is present, at 12 UTC.}
  \label{tab:mucin}
  \begin{tabular}{@{}lSS@{}}
    \toprule
    Validity time & {Present (\si{\percent})} & {Missing (\si{\percent})} \\
    \midrule
    2024-01-15 & 3.8  & 96.2 \\
    2024-02-15 & 5.7  & 94.3 \\
    2024-03-15 & 4.8  & 95.2 \\
    2024-07-15 & 28.8 & 71.2 \\
    2024-10-15 & 15.2 & 84.8 \\
    2025-02-15 & 5.5  & 94.5 \\
    \bottomrule
  \end{tabular}
\end{table}

\campo{mucape}, by contrast, is dense, with no missing points: where the routine finds no CAPE,
the published zero is the correct value. It is zero at \num{1536450} of the \num{2220273} points
(\SI{69.2}{\percent}) on 15 February 2024 and at \num{1020809} (\SI{46.0}{\percent}) on 15 July
2024, both at 12 UTC, and at \num{744745} (\SI{33.5}{\percent}) on 4 September 2024 at 15 UTC.
